\subsection{A.4. What the binomial degree filtration does and does not give}

\paragraph{A-015. Proved: the coefficient associated graded.}
Let $V$ be a free abelian module of linear coordinate functions and
$\operatorname{Num}(V)$ the ring of integer-valued rational
polynomials, filtered by total polynomial degree. Then canonically,
functorially for integral linear substitutions,
\[
 \operatorname{gr}\operatorname{Num}(V)\cong\Gamma_{\mathbb Z}(V).
 \tag{A-12}
\]
\emph{Proof.}
In coordinates $x_1,\ldots,x_d$ the products
$\prod_j\binom{x_j}{a_j}$ form an integral basis. The successive
finite-difference expansion has coefficients
$(\Delta^a f)(0)$, which are integers when $f$ is integer valued;
this proves spanning and uniqueness. Their leading terms are
$\prod_j x_j^{a_j}/a_j!$, so leading multiplication has coefficients
$\prod_j\binom{a_j+b_j}{a_j}$. This proves (A-12) and its
naturality. On the affine numerical simplex $\sum x_i=1$, eliminate
$x_0$; the leading relation is $\sum x_i=0$. Thus its coefficient
associated graded is $\Gamma(\Psi)$, whereas the universal simplex
with the variable $T=\sum x_i$ has $\Gamma(\Phi)$.
Neither ring calculation supplies a differential.\hfill$\square$

\paragraph{A-016. Refuted: the strict simplicial-DGA binomial bridge.}
There is no simplicial filtered unital differential graded
$\mathbb Z$-algebra $B$, even allowing a noncommutative product,
with the following data: $B_n^0=\operatorname{Num}(\mathbb Z^{n+1})$
with degree filtration and usual simplicial pullbacks; a central
relative coefficient algebra $\operatorname{Num}(\mathbb Z T)$
with $dT=0$; and an isomorphism of associated-graded simplicial
DGAs
\[
 \operatorname{gr}B=\Gamma(\Phi)\otimes\Lambda(\Psi)
\]
with the PD de Rham differential. The filtration is nonnegative,
exhaustive, and multiplicative. In particular $F_0B^1=0$ and
$F_1B^1=\Psi$.

\emph{Proof.}
For linear $f,g$, graded commutativity of the leading symbols gives
$[df,g]\in F_1B^1$. Constants and $T$ are central and $dT=0$,
so this defines a simplicial bilinear map
$b:\Psi\otimes\Psi\to\Psi$, $b(df,dg)=[df,g]$.
Let $\beta_2(f)=\binom f2$. The leading-symbol differential gives
$h(f)=d\beta_2(f)-f\,df\in F_1B^1$. Differentiate the exact
integer identity $2\beta_2(f)=f^2-f$. Leibniz gives
\[
 2h(f)=b(df,df)-df.
 \tag{A-13}
\]
Modulo two, therefore, $b(u,u)=u$.

Over $\mathbb F_2$ identify $\Psi_1$ with the difference of the
two vertex values and $\Psi_2$ with successive differences $(a,b)$.
The three face maps are $a$, $b$, and $a+b$. The diagonal requirement
forces $b_1(s,t)=st$. Naturality for the first two faces then forces
$b_2((a,b),(c,d))=(ac,bd)$. The third demands
$ac+bd=(a+b)(c+d)$, which fails for $(a,b)=(1,0)$ and
$(c,d)=(0,1)$. This proves the contradiction. The finite program
\texttt{A\_binomial\_obstruction.py} enumerates all $256$ bilinear
maps in dimension two; successively $16$, $1$, and $0$ survive
the three face conditions.\hfill$\square$

This refutes the specified strict interpretation of L2. The differential,
product, relative coefficients and simplicial maps are all part of the
claim. Ordinary differentiation already fails integrality
on $\binom x2$. A one-variable finite-difference calculus can repair
that using $dx\,f(x)=f(x+1)dx$, but cannot satisfy the three simplicial
face conditions above. This particular proof does not exclude
homotopy-coherent models. The later Statement~\ref{sa08:A026}
supplies the additional argument: it excludes every exhaustive
nonnegative relative filtered numerical algebra whose actual all-$X$
associated graded is the weighted PD cochain algebra, including
derived and $E_\infty$ realizations. The weaker total/shadow
comparison in Statement~\ref{sa08:D026} remains valid.

The homotopy-recovery statement proved in A.2 is independent
of this obstruction: it uses the full weighted cochains and a derived
symmetric monoidal shadow quotient, not just (A-12). A source-by-source
comparison with Horel, Toen, Mandell, and Cartan is given
in the following literature subsection; no binomial comparison theorem is assumed.

\paragraph{A-017. Proved: the ordinary numerical K\"ahler complex and its actual associated graded.}
For the one-variable image of ordinary differentiation, Kareem
\cite[Theorem~6.3.10 and Corollary~6.3.12, pp.~135--136]{A_Kareem2018}
already proves that its numerical-polynomial module span is $\Q[x]$.
We additionally identify the cotangent complex, excluding hidden
torsion, and determine the full induced degree filtration below.
Put $R_d=\operatorname{Num}(\mathbb Z^d)$ and give it the
polynomial-degree filtration. There are canonical identifications
in the derived category of $R_d$-modules
\[
 L_{R_d/\mathbb Z}\simeq
 \Omega^1_{R_d/\mathbb Z}[0]
 \simeq \mathbb Q[x_1,\ldots,x_d]\otimes_{\mathbb Z}\mathbb Z^d[0].
 \tag{A-14}
\]
Here the basis vectors on the right are $dx_j$, and the differential
of a numerical polynomial is its ordinary rational-polynomial
differential. In the universal relative simplex use the integral
linear coordinates $T,x_1,\ldots,x_n$, where
$T=\sum_{i=0}^nX_i$ and $x_j=X_j$ for $j>0$. With
$R=\operatorname{Num}(\mathbb Z^{n+1})$ and
$R_0=\operatorname{Num}(\mathbb Z T)$, the ordinary relative de Rham
complex has
\[
 \Omega^0_{R/R_0}=R,\qquad
 \Omega^q_{R/R_0}=
 \mathbb Q[T,x_1,\ldots,x_n]\otimes
 \Lambda^q\langle dx_1,\ldots,dx_n\rangle\quad(q>0).
 \tag{A-15}
\]
The induced degree filtration, rather than the filtration by rational
polynomial degree alone, is explicitly
\[
 F_w\Omega^q_{R/R_0}
 =\sum_{|a|+|b_1|+\cdots+|b_q|\leq w}
 \mathbb Z\,
 \beta_a\,d\beta_{b_1}\wedge\cdots\wedge d\beta_{b_q},
 \qquad
 \beta_a=\binom{T}{a_0}\prod_{j=1}^n\binom{x_j}{a_j}.
 \tag{A-16}
\]
For every $n,q,w$ this finite lattice presentation determines the
whole associated graded, including its torsion, its differential,
its product and its simplicial operators. In particular, for $n=1$,
\[
 F_1\Omega^1=\mathbb Z\,dx,
 \qquad
 F_2\Omega^1=
 \tfrac12\mathbb Z\,dx\ \oplus\ \mathbb Z\,x\,dx
 \ \oplus\ \mathbb Z\,T\,dx,
\]
so that
\[
 \operatorname{gr}_2\Omega^1_{R/R_0}
 \cong\mathbb Z^2\oplus\mathbb Z/2.
 \tag{A-17}
\]
The predicted PD term $\Gamma^1(\Phi_1)\otimes\Psi_1$ is free of
rank two. Thus even the underlying associated-graded module of the
ordinary numerical K\"ahler complex is not that PD complex.

\emph{Proof.}
The binomial basis proves that $R_d$ is a free abelian group. Fix a
prime $p$. Evaluation gives an injection of $R_d/pR_d$ into the
ring of functions $\mathbb Z^d\to\mathbb F_p$: an integer-valued
polynomial that takes values divisible by $p$ is $p$ times another
integer-valued polynomial. If each $a_j<p^r$, Vandermonde's identity
and $\binom{p^r}{i}\equiv0\pmod p$ for $0<i<p^r$ show that
$\beta_a$ is periodic modulo $p^r$ in every coordinate. The
evaluation matrix of $\binom xa$, $0\leq x,a<p^r$, is triangular
with diagonal entries one. Its $d$-fold tensor product is therefore
invertible over $\mathbb F_p$. Consequently these $\beta_a$ span
exactly all functions on $(\mathbb Z/p^r)^d$. Injectivity of
evaluation shows that this span is a subalgebra, and the exhaustive
binomial basis gives
\[
 R_d/pR_d=\mathop{\rm colim}_r
 \operatorname{Fun}((\mathbb Z/p^r)^d,\mathbb F_p).
 \tag{A-18}
\]
Each algebra on the right is a finite product of $\mathbb F_p$ and
is \emph{\'etale}. Its cotangent complex vanishes; the same is true
of the filtered colimit by the Stacks Project, Lemma~92.8.4,
Tag~08R2, clauses~5 and~6, which we opened.

Since $R_d$ is $\mathbb Z$-flat, derived base change
(Stacks, Lemma~92.6.2, Tag~08QQ) now gives
$L_{R_d/\mathbb Z}\otimes^{\mathbf L}_{\mathbb Z}\mathbb F_p=0$.
The tensor product is the cone of multiplication by $p$ on the
cotangent complex. Hence every prime acts invertibly on that
complex in $D(R_d)$, and all of its homology modules are
$\mathbb Q$-vector spaces. The natural rationalization map is
therefore a quasi-isomorphism. After rationalization $R_d$ is the
polynomial ring $\mathbb Q[x_1,\ldots,x_d]$, whose cotangent
complex is free on the $dx_j$ in degree zero. This proves (A-14),
including the absence of any hidden divisible torsion.

The joint binomial basis makes $R$ a free $R_0$-module on the
$x$-binomials. The transitivity triangle (Stacks,
Proposition~92.7.4, Tag~08QX) identifies $L_{R/R_0}$ with the
quotient of the rational-polynomial differential module by the
direct summand generated by $dT$. This proves the assertion in
degree one. Exterior powers over $R$ of this rationalized free
module are its ordinary rational exterior powers, giving (A-15).

By definition the induced filtered differential module is generated
by $f_0df_1\wedge\cdots\wedge df_q$ with sum of filtration degrees
at most $w$. Expand each $f_i$ in its finite binomial basis.
Multilinearity proves (A-16), with no omitted generators or
relations: it is an equality of subgroups of the rational vector
space in (A-15). In coordinates, use the finite basis
\[
 T^{c_0}x_1^{c_1}\cdots x_n^{c_n}
 dx_{i_1}\wedge\cdots\wedge dx_{i_q},
 \qquad |c|\leq w-q.
\]
Write the generators (A-16) as rational columns, clear a common
denominator, and take column Hermite normal form. Express the
previous lattice in this basis and take Smith normal form. This
gives $F_w/F_{w-1}$ as an explicit finitely presented abelian group.
Ordinary differentiation, wedge multiplication and integral linear
substitution on the same columns give all the requested structure
maps. These maps preserve the asserted filtrations by their defining
formulas, so the procedure also gives their maps on the quotients.
This is a terminating algorithm in every degree and dimension, not
only in the examples computed.

For the displayed weight-two example the nonzero generators reduce
to $dx$, $x\,dx$, $T\,dx$, and
$d\binom x2=(x-\tfrac12)dx$. They generate exactly the lattice
claimed, proving (A-17). In this associated graded,
$d[\binom x2]=[x\,dx]-[\tfrac12dx]$; the last class is the
nonzero two-torsion class. The program
\texttt{A\_numerical\_differentials.py} independently computes this
Smith quotient and eleven further quotients. All twelve finite
computations agree with the stated lattice description.
\hfill$\square$

This also explains why ordinary numerical K\"ahler forms do not
provide a de Rham resolution merely by rational differentiation.
For example $\tfrac12dx$ is a closed relative one-form for
$\operatorname{Num}(\mathbb Z[T,x])$ over
$\operatorname{Num}(\mathbb Z T)$, but has no primitive there:
a putative primitive would be $x/2+g(T)$, whose values at $x=0,1$
cannot both be integers. This failure concerns this specified
ordinary K\"ahler construction. It does not identify an unspecified
alternative filtered binomial cochain theory with it.
